\documentclass[10pt]{article}

% ---------- Document Identity (update these only) ----------
\newcommand{\DocStage}{}
\newcommand{\DocTitle}{Your Road to Tier One \textbf{SOF}}
\newcommand{\ProgramName}{182-Week Civilian-to-Selection Readiness Pipeline}
\newcommand{\ProgramSubtitle}{}

% ---------- Page ----------
\usepackage[legalpaper,left=0.5in,right=0.5in,top=0.6in,bottom=0.45in]{geometry}

% ---------- Typography ----------
\usepackage[T1]{fontenc}
\usepackage[utf8]{inputenc}
\usepackage{libertinus}
\usepackage{microtype}
\usepackage{parskip}
\usepackage{ragged2e}
\usepackage{textcomp}
\AtBeginDocument{\color{black}}
\AtBeginDocument{\pagecolor{white}}
\AtBeginDocument{\justifying}

% ---------- Landscape pages ----------
\usepackage{pdflscape}

% ---------- Color ----------
\usepackage[table]{xcolor}
\definecolor{StageBlue}{HTML}{1F4E79}
\definecolor{StageBlueDark}{HTML}{163A5A}
\definecolor{LightGray}{HTML}{F2F4F7}
\definecolor{MidGray}{HTML}{D9DEE5}
\definecolor{RuleGray}{HTML}{C7CED8}
\definecolor{HeaderInk}{HTML}{000000}
\definecolor{Ink}{HTML}{000000}

% ---------- Utility ----------
\usepackage{enumitem}
\sloppy
\setlength{\emergencystretch}{2em}

% ---------- Boxes / UI ----------
\usepackage[most]{tcolorbox}

\tcbset{
  charterTitle/.style={
    enhanced,
    colback=StageBlue!4,
    colframe=StageBlue,
    boxrule=1.25pt,
    arc=3pt,
    left=16pt,right=16pt,top=14pt,bottom=14pt,
    borderline north={3.2pt}{0pt}{StageBlue},
    borderline west={4pt}{0pt}{StageBlue},
    borderline south={1.25pt}{0pt}{StageBlue},
    drop shadow={black!25!white},
  },
  charterRule/.style={
    enhanced,
    colback=LightGray,
    colframe=RuleGray,
    boxrule=0.85pt,
    arc=3pt,
    left=12pt,right=12pt,top=10pt,bottom=10pt,
    borderline west={3pt}{0pt}{StageBlue},
  },
  charterBadge/.style={
    enhanced,
    colback=StageBlue,
    coltext=white,
    colframe=StageBlueDark,
    boxrule=0.6pt,
    arc=2pt,
    left=6pt,right=6pt,top=3pt,bottom=3pt,
  }
}
\newtcbox{\Badge}{nobeforeafter,charterBadge}

% ---------- Lists ----------
\setlist[itemize]{leftmargin=1.5em, itemsep=0.25em, topsep=0.25em}
\setlist[enumerate]{leftmargin=1.8em, itemsep=0.25em, topsep=0.25em}

% ---------- TOC ----------
\usepackage{tocloft}
\setcounter{tocdepth}{3}
\setlength{\cftbeforesecskip}{3pt}
\setlength{\cftbeforesubsecskip}{2pt}
\setlength{\cftbeforesubsubsecskip}{0pt}
\renewcommand{\cftsecfont}{\bfseries}
\renewcommand{\cftsecpagefont}{\bfseries}
\renewcommand{\cftsubsecfont}{\normalfont}
\renewcommand{\cftsubsecpagefont}{\normalfont}
\renewcommand{\cftsubsubsecfont}{\normalfont}
\renewcommand{\cftsubsubsecpagefont}{\normalfont}
\setlength{\cftsecindent}{0pt}
\setlength{\cftsubsecindent}{1.25em}
\setlength{\cftsubsubsecindent}{2.8em}
\setlength{\cftsecnumwidth}{2.1em}
\setlength{\cftsubsecnumwidth}{2.8em}
\setlength{\cftsubsubsecnumwidth}{3.6em}

% Remove the built-in TOC heading so only the custom heading is shown
\makeatletter
\renewcommand{\tableofcontents}{\@starttoc{toc}}
\makeatother

% ---------- Header/Footer: page number only ----------
\usepackage{fancyhdr}
\pagestyle{fancy}
\fancyhf{}
\renewcommand{\headrulewidth}{0pt}
\renewcommand{\footrulewidth}{0pt}
\fancyfoot[C]{\thepage}

% ---------- Section formatting ----------
\usepackage{titlesec}

\titleformat{\section}
  {\Large\bfseries\color{Ink}}
  {\thesection}{0.6em}{}
  [\vspace{0.25em}\textcolor{StageBlue}{\rule{\linewidth}{0.8pt}}]

\titleformat{\subsection}
  {\large\bfseries\color{Ink}}
  {\thesubsection}{0.6em}{}

\titleformat{\subsubsection}
  {\normalsize\bfseries\color{Ink}}
  {\thesubsubsection}{0.6em}{}

\titlespacing*{\section}{0pt}{0.95em}{0.35em}
\titlespacing*{\subsection}{0pt}{0.65em}{0.25em}
\titlespacing*{\subsubsection}{0pt}{0.55em}{0.2em}

% ---------- Table engine ----------
\usepackage{tabularray}
\UseTblrLibrary{booktabs}

% ---------- tabularray: GLOBAL table treatment ----------
\SetTblrInner{
  cells = {fg=black, bg=white, valign=t},
}

% ---------- Header helpers ----------
\newcommand{\Hdr}[1]{\shortstack[c]{\strut #1\strut}}

% ---------- Lines ----------
\newcommand{\TitleLine}{\textcolor{StageBlue}{\rule{0.98\linewidth}{0.95pt}}}
\newcommand{\SubTitleLine}{\textcolor{RuleGray}{\rule{0.98\linewidth}{0.65pt}}}

% ---------- Continued on next page ----------
\newcommand{\ContinuedOnNextPageInline}{%
  \par
  \vfill
  \noindent\textcolor{RuleGray}{\rule{\linewidth}{1pt}}\vspace{-2pt}
  \noindent\hfill\textit{Continued on next page}\par\vspace{-10pt}
  \noindent\textcolor{RuleGray}{\rule{\linewidth}{1pt}}
  \clearpage
}

% ---------- PDF hyperlinks ----------
\usepackage[hidelinks]{hyperref}
\hypersetup{
  colorlinks=false,
  pdfborder={0 0 0},
  linktoc=all,
  pdftitle={\DocTitle},
  pdfauthor={\ProgramName}
}

% =========================================================
\begin{document}

\section*{10.D.1 — Ladder Framework — How Ladders Work}
\phantomsection
\addcontentsline{toc}{section}{10.D.1 — Ladder Framework — How Ladders Work}

\begin{itemize}
\item Verify data first:
  \begin{itemize}
  \item trend-based posture only; no single-day decisions.
  \item use Stage 2 metrics and Stage 2.E logging posture as the measurement interface.
  \item weekly bodyweight average and weekly circumference suite are primary \textbf{BODYCOMP} signals; recovery markers and session RPE are constraint signals.
  \end{itemize}
\item Fix compliance before changing macros:
  \begin{itemize}
  \item apply 10.B lever order: simplify and restore compliance before calorie math.
  \item compliance thresholds are gating conditions:
    \begin{itemize}
    \item $\geq$6/7 on-plan days for minor tweaks,
    \item $\geq$12/14 on-plan days for structural calorie changes.
    \end{itemize}
  \end{itemize}
\item Enforce hold periods:
  \begin{itemize}
  \item 7–14 days after macro changes unless safety requires immediate reversal.
  \item maintenance reset is 7 days at maintenance when triggered (10.B).
  \end{itemize}
\item No novelty near gates and test weeks:
  \begin{itemize}
  \item deload/test weeks: calories set to maintenance (10.B).
  \item test-day carb redistribution allowed per 10.B (+50–100 g carbs offset by fats).
  \item novelty foods/products prohibited near gates (final 7–10 days).
  \item caffeine routine changes are novelty: no caffeine pattern changes within 14 days of gate windows.
  \end{itemize}
\item Program-wide Do Not Do set:
  \begin{itemize}
  \item do not crash diet or attempt \textgreater0.5\% \textbf{BW}/week planned average loss.
  \item do not add stimulants to force adherence; do not use electrolytes as stimulant substitutes (see 10.C; cross-reference 7.05 boundaries by \textbf{ID} only).
  \item do not add multiple new changes in the same week.
  \item do not attempt to make up missed compliance with extra restriction, dehydration, or “compensation days.”
  \item do not force test attempts when heat/illness/safety conditions invalidate validity; reslot per Stage 3.
  \end{itemize}
\end{itemize}

\newpage

Required ladder template table:

\begingroup
\scriptsize
\setlength{\tabcolsep}{2pt}
\renewcommand{\arraystretch}{1.12}

\begin{longtblr}{
  width=\linewidth,
  rowhead=1,
  colspec = {
    X[l,3.20]
    X[l,1.65]
    X[l,2.40]
    X[l,2.40]
    X[l,2.35]
  },
  hlines,
  vlines,
  cells = {hsep=6pt, vsep=6pt, halign=l, valign=t},
  row{odd} = {bg=white},
  row{even} = {bg=LightGray},
  row{1} = {bg=StageBlue!15, fg=black, font=\bfseries, halign=l, valign=m},
}
\Hdr{Ladder Name} &
\Hdr{Primary Signal} &
\Hdr{First Check} &
\Hdr{Escalation Trigger} &
\Hdr{Stop Condition} \\
Plateau Ladder — Bodyweight Trend Stalled &
14-day trend stall &
Compliance $\geq$12/14 and waist method consistency &
Persistent stall after one step + hold &
Red-flag symptoms; restrictive behavior risk \\
Too-Fast Loss Ladder — Recovery/Performance Risk &
\textgreater0.75\% \textbf{BW}/week $\times$ 2 weeks &
Confirm weekly average + compliance $\geq$12/14 &
Sleep/performance collapse &
Red-flag symptoms \\
Unexpected Gain Ladder — Trend Up &
2-week upward trend &
Sodium/carb swings + compliance &
Sustained upward drift under high compliance &
Red-flag symptoms \\
Hunger/Cravings Ladder — Adherence Breakdown Risk &
cravings/urge-to-break plan &
Compliance trend + sleep trend &
Repeated off-plan days &
Restrictive behavior escalation \\
Low Energy/High RPE Ladder — Recovery Collapse Risk &
RPE trend up; fatigue up &
Sleep markers + illness flags &
\textbf{MODIFIED}/\textbf{INVALID} sessions due to fatigue &
Red-flag symptoms \\
Poor Sleep Trend Ladder &
sleep hours/quality down &
Caffeine timing stability + stressors &
Persistent multi-week degradation &
Panic/dissociation; self-harm statements \\
GI Distress Ladder — Training or Daily &
GI distress persistent &
Novelty foods/products check &
Persistent or severe symptoms &
Red-flag symptoms; dehydration risk \\
Cramping/Dehydration Symptoms Ladder — 10.C Linkage &
cramps/dizziness &
Heat/cold + fluid/electrolyte logging &
Recurrent symptoms &
Heat illness suspicion; syncope \\
Test-Week Performance Dip Ladder — Validity Protection &
dip in test-week output &
No novelty confirmed; maintenance week applied &
Repeat dips with symptoms &
Illness/heat triggers invalidate testing \\
Budget Constraint/Food Access Instability Ladder &
access instability &
protein adequacy + baseline calories &
Repeated under-eating &
Restrictive behavior risk \\
Travel/Schedule Disruption Ladder &
disrupted schedule &
portable baseline + hydration continuity &
Repeated missed logs/sessions &
Red-flag symptoms \\
Water Retention Perception Ladder &
scale fluctuations &
7–14 day trend check &
persistent upward trend &
Red-flag symptoms \\
\end{longtblr}

\endgroup

\newpage

\section*{10.D.2 — Troubleshooting Ladders — Required Set}
\phantomsection
\addcontentsline{toc}{section}{10.D.2 — Troubleshooting Ladders — Required Set}

\subsection*{LADDER-01 — Plateau Ladder — Bodyweight Trend Stalled}
\phantomsection
\addcontentsline{toc}{subsection}{LADDER-01 — Plateau Ladder — Bodyweight Trend Stalled}

Primary Signal: 14 days with \textless0.10\% bodyweight change and no meaningful waist reduction.\\
Confirm With Data: Weekly bodyweight average (\textbf{MET-2B-001}) and weekly waist circumference (\textbf{MET-2B-007}) with consistent landmark/tape posture.\\
Stop Conditions Safety: Red-flag symptoms; severe fatigue collapse; restrictive behavior escalation signs.\\
Review Cadence: Weekly.\\
Do Not Do: Crash diet; stack multiple changes; cut during deload/test weeks; introduce novelty near gates.\\
Steps in order:

\begin{itemize}
\item Step 0 — Verify signal quality
  \begin{itemize}
  \item Trend window: last 14 days.
  \item Compliance check: structural change requires $\geq$12/14 on-plan days. If \textless12/14, plateau is not confirmed for action; treat as compliance problem.
  \item Waist validity: confirm consistent tape method and landmark. If inconsistent, treat waist as non-comparable and use bodyweight trend only.
  \end{itemize}
\item Step 1 — Lowest-cost correction
  \begin{itemize}
  \item Compliance simplification: remove optional complexity; standardize staples; reduce discretionary snacks; maintain protein and fat floor per 10.A/10.B.
  \item Confirm no novelty events occurred in the last 7–10 days if gates are near.
  \end{itemize}
\item Step 2 — Next correction
  \begin{itemize}
  \item Tighten discretionary carbohydrates first (within current calorie posture) while keeping protein stable and fat floor 20–30\% intact.
  \end{itemize}
\item Step 3 — Macro adjustment
  \begin{itemize}
  \item If compliance $\geq$12/14 and plateau persists: reduce calories by one step (-150 to -250 kcal/day) per 10.B; hold 7–14 days.
  \end{itemize}
\item Step 4 — Escalation, hold, refer
  \begin{itemize}
  \item If plateau persists after one step + full hold period: execute a second compliance audit and only then consider a second single-step reduction (do not stack in one week).
  \item If gates/test weeks are within 7–10 days: defer structural change; protect validity; reslot if needed per Stage 3.
  \item Hold Period Before Next Change: 7–14 days after Step 3.
  \item What to Log: Weekly \textbf{BW} average, waist method consistency, on-plan count, change date, step size, hold period, gate proximity Y/N.
  \item Worked Example:
    \begin{itemize}
    \item Step 0: 14-day \textbf{BW} average change = 0.05\%; waist measures unchanged; compliance 13/14; waist method stable.
    \item Step 1: simplify snacks; standardize breakfast/lunch staples for 7 days.
    \item Step 2: reduce discretionary carbs (e.g., remove one routine carb add-on) without changing protein.
    \item Step 3: after continued plateau, reduce calories by 200 kcal/day and hold 10 days.
    \item Step 4: if still stalled and no gate window imminent, run a second compliance audit before any additional step.
    \end{itemize}
  \end{itemize}
\end{itemize}

\newpage

\subsection*{LADDER-02 — Too-Fast Loss Ladder — Recovery/Performance Risk}
\phantomsection
\addcontentsline{toc}{subsection}{LADDER-02 — Too-Fast Loss Ladder — Recovery/Performance Risk}

Primary Signal: \textgreater0.75\% bodyweight per week sustained for 2 consecutive weeks.\\
Confirm With Data: Weekly \textbf{BW} average trend (\textbf{MET-2B-001}) across two weeks + compliance $\geq$12/14.\\
Stop Conditions Safety: Red-flag symptoms; severe dizziness; near-syncope; severe GI distress.\\
Review Cadence: Weekly.\\
Do Not Do: Keep cutting “because it’s working”; add stimulants; restrict fluids.\\
Steps in order:

\begin{itemize}
\item Step 0 — Verify signal quality
  \begin{itemize}
  \item Trend window: last 14 days.
  \item Compliance: $\geq$12/14 required to treat as true signal rather than noise.
  \item Check for acute water shifts (high sodium/carb swings) before acting.
  \end{itemize}
\item Step 1 — Lowest-cost correction
  \begin{itemize}
  \item Ensure deload/test posture is not being violated (maintenance week rules).
  \item Confirm hydration baseline and session top-ups align to 10.C; do not interpret dehydration as “fat loss.”
  \end{itemize}
\item Step 2 — Next correction
  \begin{itemize}
  \item Shift distribution: increase carbohydrates modestly while keeping protein stable; keep fat within 20–30\%.
  \end{itemize}
\item Step 3 — Macro adjustment
  \begin{itemize}
  \item Raise calories by one step (+150 to +250 kcal/day) per 10.B; preferentially raise carbs while keeping fat floor intact.
  \end{itemize}
\item Step 4 — Escalation, hold, refer
  \begin{itemize}
  \item If fast loss continues with recovery decline: execute 7-day maintenance reset per 10.B; resume prior posture only when sleep/fatigue stabilize.
  \item If red-flag symptoms appear: stop and seek evaluation.
  \item Hold Period Before Next Change: 7–14 days after Step 3 unless safety requires maintenance reset sooner.
  \item What to Log: Two-week \textbf{BW} rate, sleep/fatigue markers, calorie step change, carb distribution change, hold period.
  \item Worked Example:
    \begin{itemize}
    \item Step 0: Week 1 loss 0.9\%, Week 2 loss 0.8\%, compliance 14/14.
    \item Step 2: increase carbs slightly (no novelty foods).
    \item Step 3: +200 kcal/day; hold 10 days.
    \item Step 4: if fatigue remains elevated, apply 7-day maintenance reset.
    \end{itemize}
  \end{itemize}
\end{itemize}

\newpage

\subsection*{LADDER-03 — Unexpected Gain Ladder — Trend Up}
\phantomsection
\addcontentsline{toc}{subsection}{LADDER-03 — Unexpected Gain Ladder — Trend Up}

Primary Signal: Perceived gain or a 2-week upward trend in weekly \textbf{BW} average.\\
Confirm With Data: Weekly \textbf{BW} averages across 14 days + waist trend only if method consistent; compliance thresholds apply.\\
Stop Conditions Safety: Red-flag symptoms; severe GI distress; suspected illness with systemic symptoms.\\
Review Cadence: Weekly.\\
Do Not Do: Cut immediately based on a single weigh-in; add dehydration tactics.\\
Steps in order:

\begin{itemize}
\item Step 0 — Verify signal quality
  \begin{itemize}
  \item Trend window: 7–14 days; primary decision uses 14-day trend.
  \item Compliance check: if \textless12/14, treat as compliance problem, not a metabolism problem.
  \end{itemize}
\item Step 1 — Lowest-cost correction
  \begin{itemize}
  \item Check for sodium/carb swing and novelty events (new foods, new eating window, new caffeine pattern).
  \item Re-stabilize to baseline staples and routine.
  \end{itemize}
\item Step 2 — Next correction
  \begin{itemize}
  \item Execute a compliance rebuild week (target $\geq$6/7 on-plan) with simplified, repeatable choices.
  \end{itemize}
\item Step 3 — Macro adjustment
  \begin{itemize}
  \item Only if a 14-day upward trend persists under $\geq$12/14 compliance: reduce calories by one step (-150 to -250) and hold 7–14 days.
  \end{itemize}
\item Step 4 — Escalation, hold, refer
  \begin{itemize}
  \item If trend continues upward and recovery markers degrade, hold and reassess; do not stack cuts; consider maintenance reset if diet fatigue appears.
  \item Hold Period Before Next Change: 7–14 days after Step 3.
  \item What to Log: 2-week trend, compliance counts, novelty events, step change and hold.
  \item Worked Example:
    \begin{itemize}
    \item Step 0: single-day +3 lb after high sodium/carb day; weekly average unchanged.
    \item Step 1: stabilize food choices for 7 days; no macro changes.
    \item Step 3: not applied because no 2-week trend confirmed.
    \end{itemize}
  \end{itemize}
\end{itemize}

\newpage

\subsection*{LADDER-04 — Hunger/Cravings Ladder — Adherence Breakdown Risk}
\phantomsection
\addcontentsline{toc}{subsection}{LADDER-04 — Hunger/Cravings Ladder — Adherence Breakdown Risk}

Primary Signal: escalating hunger/cravings leading to repeated off-plan days.\\
Confirm With Data: on-plan day count ($\geq$6/7 minor threshold), sleep quality trend, fatigue trend, and whether deficit is too aggressive.\\
Stop Conditions Safety: restrictive behavior escalation signs; binge/purge behaviors; fear-driven avoidance of eating.\\
Review Cadence: Weekly, with mid-week check-in if adherence is collapsing.\\
Do Not Do: punitive restriction; “compensation days”; stimulant escalation.\\
Steps in order:

\begin{itemize}
\item Step 0 — Verify signal quality
  \begin{itemize}
  \item Trend window: last 7 days for adherence; last 14 days for weight trend.
  \item Compliance threshold check: if \textless6/7 on-plan, treat as system failure risk; macro changes are deferred until simplification is executed.
  \end{itemize}
\item Step 1 — Lowest-cost correction
  \begin{itemize}
  \item Increase structure, not restriction: standardize meal timing; reduce decision points; maintain protein target.
  \item Add “planned consistency” options rather than reactive snacking (no meal plan generated; just simplification posture).
  \end{itemize}
\item Step 2 — Next correction
  \begin{itemize}
  \item Rebalance within current calories: increase volume foods and fiber posture using familiar foods; avoid novelty high-fiber changes near gates.
  \end{itemize}
\item Step 3 — Macro adjustment
  \begin{itemize}
  \item If cravings persist and compliance is $\geq$12/14 but hunger remains high: increase calories by +150 kcal/day (prefer carbs) OR reduce deficit by one step, hold 7–14 days.
  \end{itemize}
\item Step 4 — Escalation, hold, refer
  \begin{itemize}
  \item If restrictive behavior warning signs appear: maintenance reset 7 days; seek professional support posture; do not intensify restriction.
  \item Hold Period Before Next Change: 7–14 days after Step 3; immediate maintenance reset if restrictive behavior risk is detected.
  \item What to Log: hunger/craving notes, on-plan count, sleep trend, change executed, hold period, gate proximity Y/N.
  \item Worked Example:
    \begin{itemize}
    \item Step 0: on-plan days 4/7; cravings high; sleep poor.
    \item Step 1: simplify to repeatable meals and consistent meal timing for 7 days.
    \item Step 3: not allowed until compliance improves; once compliance reaches 12/14 and hunger persists, +150 kcal/day.
    \end{itemize}
  \end{itemize}
\end{itemize}

\newpage

\subsection*{LADDER-05 — Low Energy / High RPE Ladder — Recovery Collapse Risk}
\phantomsection
\addcontentsline{toc}{subsection}{LADDER-05 — Low Energy / High RPE Ladder — Recovery Collapse Risk}

Primary Signal: persistent low energy and rising session RPE trend despite stable training plan.\\
Confirm With Data: \textbf{MET-2B-042} (session RPE), \textbf{MET-2B-041} fatigue, \textbf{MET-2B-038} sleep hours, \textbf{MET-2B-039} sleep quality; confirm illness flags.\\
Stop Conditions Safety: red-flag symptoms; inability to function; suspected heat illness; severe dizziness.\\
Review Cadence: Weekly, with immediate action if safety signals occur.\\
Do Not Do: cut calories further; add stimulants; “push through” with more restriction.\\
Steps in order:

\begin{itemize}
\item Step 0 — Verify signal quality
  \begin{itemize}
  \item Trend window: 7 days for recovery markers and RPE; 14 days for weight trend.
  \item Compliance threshold: macro reductions are prohibited when recovery collapse is present; do-not-cut posture applies even if weight loss is “too slow.”
  \end{itemize}
\item Step 1 — Lowest-cost correction
  \begin{itemize}
  \item Stabilize sleep opportunity and hydration posture (10.C baseline + session top-ups).
  \item Simplify nutrition execution to reduce friction and missed intake.
  \end{itemize}
\item Step 2 — Next correction
  \begin{itemize}
  \item Redistribution within calories: shift carbs toward training days while keeping protein and fat floor intact; no novelty foods.
  \end{itemize}
\item Step 3 — Macro adjustment
  \begin{itemize}
  \item Execute maintenance reset: 7 days at maintenance per 10.B; exit only when sleep/fatigue stabilize.
  \end{itemize}
\item Step 4 — Escalation, hold, refer
  \begin{itemize}
  \item If recovery markers do not stabilize after maintenance reset or red flags occur: stop escalation and seek evaluation posture; reslot tests per Stage 3.
  \item Hold Period Before Next Change: maintenance reset 7 days; then 7-day reassess; no structural deficit changes until stable.
  \item What to Log: RPE trend, fatigue/sleep, hydration notes, maintenance reset start/end, review date.
  \item Worked Example:
    \begin{itemize}
    \item Step 0: RPE elevated for 6 sessions; fatigue 8/10; sleep 5 hrs/night for 4 nights.
    \item Step 3: apply 7-day maintenance reset; hold all cuts; keep protein stable.
    \item Step 4: if still collapsing, reslot any tests and seek evaluation posture.
    \end{itemize}
  \end{itemize}
\end{itemize}

\newpage

\subsection*{LADDER-06 — Poor Sleep Trend Ladder}
\phantomsection
\addcontentsline{toc}{subsection}{LADDER-06 — Poor Sleep Trend Ladder}

Primary Signal: sleep hours and/or sleep quality trending down with performance friction.\\
Confirm With Data: \textbf{MET-2B-038} sleep hours; \textbf{MET-2B-039} sleep quality; \textbf{MET-2B-041} fatigue; session validity distribution for fatigue-driven \textbf{MODIFIED}/\textbf{INVALID}.\\
Stop Conditions Safety: panic/dissociation; self-harm statements; inability to self-regulate; red-flag symptoms.\\
Review Cadence: Weekly; immediate escalation if \textbf{STOP} \textbf{NOW} triggers appear.\\
Do Not Do: cut calories further; introduce new stimulants; change caffeine pattern within 14 days of gates.\\
Steps in order:

\begin{itemize}
\item Step 0 — Verify signal quality
  \begin{itemize}
  \item Trend window: 7 days sleep markers; check compliance thresholds for any nutrition change.
  \item Confirm no caffeine pattern changes occurred within 14 days of upcoming gate windows.
  \end{itemize}
\item Step 1 — Lowest-cost correction
  \begin{itemize}
  \item Compliance simplification to reduce late-night friction: stabilize meal timing; avoid late large boluses; maintain hydration distribution rule (10.C).
  \end{itemize}
\item Step 2 — Next correction
  \begin{itemize}
  \item Hold deficit; shift toward maintenance if sleep debt persists; prioritize routine stability.
  \end{itemize}
\item Step 3 — Macro adjustment
  \begin{itemize}
  \item Maintenance reset 7 days if sleep and fatigue are collapsing (10.B); keep protein stable.
  \end{itemize}
\item Step 4 — Escalation, hold, refer
  \begin{itemize}
  \item If severe/persistent sleep disruption or mental \textbf{STOP} \textbf{NOW} triggers occur: seek professional support posture; reslot tests per Stage 3; do not intensify restriction.
  \item Optional category reference only: 7.06 (Recovery \& Sleep Aids) as non-required support posture; no dosing.
  \item Hold Period Before Next Change: 7 days after reset; then reassess; no structural deficit changes until stable.
  \item What to Log: sleep markers, caffeine routine stability, nutrition posture (hold/reset), gate proximity Y/N.
  \item Worked Example:
    \begin{itemize}
    \item Step 0: sleep hours \textless6 for 5 of 7 nights; fatigue high; no gate within 14 days.
    \item Step 3: execute maintenance reset 7 days; no caffeine pattern changes.
    \item Step 4: if panic/dissociation occurs, \textbf{STOP} \textbf{NOW} and seek professional support.
    \end{itemize}
  \end{itemize}
\end{itemize}

\newpage

\subsection*{LADDER-07 — GI Distress Ladder — Training or Daily}
\phantomsection
\addcontentsline{toc}{subsection}{LADDER-07 — GI Distress Ladder — Training or Daily}

Primary Signal: nausea, diarrhea, bloating, reflux, or significant GI discomfort interfering with training or daily function.\\
Confirm With Data: symptom notes + novelty check (new foods, fiber load changes, new electrolyte concentration, new caffeine pattern); hydration notes (10.C); illness flags.\\
Stop Conditions Safety: severe GI distress, persistent vomiting/diarrhea, dizziness/near-syncope, dehydration symptoms.\\
Review Cadence: Weekly; immediate action if severe.\\
Do Not Do: introduce new “fix” foods; add multiple changes; change fiber dramatically near gates.\\
Steps in order:

\begin{itemize}
\item Step 0 — Verify signal quality
  \begin{itemize}
  \item Trend window: last 72 hours for acute; last 7 days for pattern.
  \item Compliance thresholds still apply for structural macro changes; GI stability is a safety constraint.
  \end{itemize}
\item Step 1 — Lowest-cost correction
  \begin{itemize}
  \item Remove novelty: revert to familiar GI-safe staples; stabilize fiber to the prior known-tolerated level; hold caffeine routine (no changes near gates).
  \end{itemize}
\item Step 2 — Next correction
  \begin{itemize}
  \item Reduce high-risk items for GI tolerance (education-only posture): excessive fat boluses, large late meals, untested sugar alcohols; maintain protein target without adding new products.
  \end{itemize}
\item Step 3 — Macro adjustment
  \begin{itemize}
  \item If GI distress is driven by volume or timing and compliance is high: redistribute calories across the day without changing totals; if severe and persistent, hold deficit and consider maintenance reset.
  \end{itemize}
\item Step 4 — Escalation, hold, refer
  \begin{itemize}
  \item Severe/persistent symptoms → stop and seek evaluation posture; reslot tests per Stage 3; do not attempt gate tests under GI distress.
  \item Hold Period Before Next Change: 7 days after any structural posture change; immediate escalation if severe.
  \item What to Log: symptom type/severity, novelty check, hydration/electrolyte use, actions taken, gate proximity Y/N.
  \item Worked Example:
    \begin{itemize}
    \item Step 0: GI distress began after new high-fiber food 5 days before a planned test.
    \item Step 1: revert to prior staples; hold; no new products.
    \item Step 4: if symptoms persist, reslot test and seek evaluation posture.
    \end{itemize}
  \end{itemize}
\end{itemize}

\newpage

\subsection*{LADDER-08 — Cramping/Dehydration Symptoms Ladder — 10.C Linkage}
\phantomsection
\addcontentsline{toc}{subsection}{LADDER-08 — Cramping/Dehydration Symptoms Ladder — 10.C Linkage}

Primary Signal: cramping, dizziness, headache, unusual fatigue, or thirst issues during/after sessions.\\
Confirm With Data: 10.C intake posture (baseline + session top-ups), environment (heat/cold), electrolyte trigger conditions, symptom logs.\\
Stop Conditions Safety: confusion, fainting/near-fainting, chest pain/pressure, severe dizziness, suspected heat illness, persistent vomiting/diarrhea.\\
Review Cadence: Weekly; immediate action for safety triggers.\\
Do Not Do: chugging; exceed 32 oz/hr; new electrolyte product or concentration near gates; stimulant-combined electrolyte products.\\
Steps in order:

\begin{itemize}
\item Step 0 — Verify signal quality
  \begin{itemize}
  \item Trend window: last 7 days; check whether symptoms are isolated or recurring.
  \item Confirm environment category and whether per-hour fluid bands were respected; confirm no forced drinking and not exceeding 32 oz/hr.
  \end{itemize}
\item Step 1 — Lowest-cost correction
  \begin{itemize}
  \item Apply 10.C baseline hydration distribution rule and session sip model within 8–16 oz/hr cool/indoor or 16–24 oz/hr hot/outdoor.
  \end{itemize}
\item Step 2 — Next correction
  \begin{itemize}
  \item If $\geq$90 min duration or heavy sweat/heat or ruck: activate electrolyte posture per 10.C bands (do not invent new targets).
  \item Maintain stable product type and mixing posture; no novelty near gates/test weeks.
  \end{itemize}
\item Step 3 — Macro adjustment
  \begin{itemize}
  \item If repeated symptoms coincide with energy deficit and recovery collapse: apply 10.B maintenance reset 7 days; do not cut further.
  \end{itemize}
\item Step 4 — Escalation, hold, refer
  \begin{itemize}
  \item If heat triggers are exceeded (10.C.5): reslot tests; substitute indoors; do not test.
  \item If red-flag symptoms occur: stop and seek evaluation posture.
  \item Hold Period Before Next Change: 7 days after any maintenance reset; electrolyte changes held stable through protected windows.
  \item What to Log: environment category, fluid estimates, electrolytes Y/N, symptom flags, any reslot decisions, gate proximity Y/N.
  \item Worked Example:
    \begin{itemize}
    \item Step 0: cramps occurred on a 2-hour hot outdoor session; intake was minimal; no electrolytes used.
    \item Step 1: apply 16–24 oz/hr sip model next similar session.
    \item Step 2: activate electrolytes per 10.C trigger; maintain stable mixing.
    \item Step 4: if heat index exceeds 95\textdegree F, reslot and move indoors.
    \end{itemize}
  \end{itemize}
\end{itemize}

\newpage

\subsection*{LADDER-09 — Test-Week Performance Dip Ladder — Validity Protection}
\phantomsection
\addcontentsline{toc}{subsection}{LADDER-09 — Test-Week Performance Dip Ladder — Validity Protection}

Primary Signal: test-week output dips or perceived underperformance during protected windows.\\
Confirm With Data: confirm validity conditions were met (Stage 2.E logging complete; no novelty); confirm deload/test nutrition posture applied (10.B maintenance week rule).\\
Stop Conditions Safety: illness, heat triggers exceeded, red-flag symptoms; validity compromised → do not use for gate credit.\\
Review Cadence: Weekly; immediate reslot decision when validity is compromised.\\
Do Not Do: cut calories during test week; add new supplements; change caffeine pattern within 14 days.\\
Steps in order:

\begin{itemize}
\item Step 0 — Verify signal quality
  \begin{itemize}
  \item Confirm test was \textbf{VALID} and logged correctly; if \textbf{INVALID} or missing required fields, treat as not yet tested and reslot per Stage 3.
  \item Confirm no novelty foods/products were introduced in the final 7–10 days pre-gate.
  \end{itemize}
\item Step 1 — Lowest-cost correction
  \begin{itemize}
  \item Apply test-week posture: calories at maintenance; protein stable; no novelty; hydration stable per 10.C.
  \end{itemize}
\item Step 2 — Next correction
  \begin{itemize}
  \item If test day: apply 10.B test-day carb redistribution (+50–100 g carbs offset by fats) only if already within allowed window and using familiar foods.
  \end{itemize}
\item Step 3 — Macro adjustment
  \begin{itemize}
  \item Do not adjust macros for a single test week dip; hold 7–14 days and reassess trend and recovery markers.
  \item If repeated dips with recovery collapse: maintenance reset 7 days (10.B).
  \end{itemize}
\item Step 4 — Escalation, hold, refer
  \begin{itemize}
  \item Reslot tests when heat/illness/safety conditions invalidate validity per Stage 3; do not force.
  \item Hold Period Before Next Change: 7–14 days; no structural changes within protected windows.
  \item What to Log: validity tag, conditions notes, novelty check, maintenance week confirmation, test-day carb offset used Y/N.
  \item Worked Example:
    \begin{itemize}
    \item Step 0: performance dip occurred but Pool \textbf{ID} missing; test is inadmissible → reslot, no macro change.
    \item Step 1: ensure maintenance week posture next attempt; keep diet stable.
    \item Step 2: apply carb offset on test day if already established and GI-safe.
    \end{itemize}
  \end{itemize}
\end{itemize}

\newpage

\subsection*{LADDER-10 — Budget Constraint / Food Access Instability Ladder}
\phantomsection
\addcontentsline{toc}{subsection}{LADDER-10 — Budget Constraint / Food Access Instability Ladder}

Primary Signal: inability to obtain or prepare planned foods consistently; repeated off-plan days due to access constraints.\\
Confirm With Data: on-plan day counts; whether misses are access-driven; weight trend and recovery markers.\\
Stop Conditions Safety: restrictive behavior escalation; severe under-eating due to access.\\
Review Cadence: Weekly.\\
Do Not Do: treat budget constraints as justification for starvation; do not cut further when access is unstable.\\
Steps in order:

\begin{itemize}
\item Step 0 — Verify signal quality
  \begin{itemize}
  \item Trend window: last 7 days for compliance; 14 days for weight trend.
  \item Compliance thresholds apply: structural changes require $\geq$12/14; if not met, simplify first.
  \end{itemize}
\item Step 1 — Lowest-cost correction
  \begin{itemize}
  \item Switch to staple-based minimal execution posture (no meal plans): simple, repeatable foods; consistent protein anchor; stable hydration.
  \end{itemize}
\item Step 2 — Next correction
  \begin{itemize}
  \item Reduce variety and reduce “special” items; prioritize the smallest set of repeatable foods that hit protein and calorie posture.
  \end{itemize}
\item Step 3 — Macro adjustment
  \begin{itemize}
  \item If access instability prevents consistent deficit execution, move to maintenance temporarily; hold 7 days, then reassess.
  \end{itemize}
\item Step 4 — Escalation, hold, refer
  \begin{itemize}
  \item If restrictive behavior risk increases or under-eating becomes persistent: maintenance reset and professional support posture; do not intensify restriction.
  \item Hold Period Before Next Change: 7 days at maintenance if executed; then weekly reassess.
  \item What to Log: access constraint, simplified staple posture adopted, compliance counts, any maintenance reset.
  \item Worked Example:
    \begin{itemize}
    \item Step 0: on-plan 3/7 due to food access; weight trend unclear.
    \item Step 1: adopt staple-based baseline and stabilize for 7 days.
    \item Step 3: hold maintenance until access stabilizes.
    \end{itemize}
  \end{itemize}
\end{itemize}

\newpage

\subsection*{LADDER-11 — Travel / Schedule Disruption Ladder}
\phantomsection
\addcontentsline{toc}{subsection}{LADDER-11 — Travel / Schedule Disruption Ladder}

Primary Signal: travel/work disruption causing missed tracking, irregular meals, and hydration drift.\\
Confirm With Data: on-plan day count, logging completeness (Stage 2.E posture), hydration fields, sleep markers.\\
Stop Conditions Safety: red-flag symptoms; severe sleep deprivation; illness flags.\\
Review Cadence: Weekly; travel week treated as “stability week.”\\
Do Not Do: attempt aggressive deficit changes while travel disrupts compliance; change caffeine pattern within 14 days of gates.\\
Steps in order:

\begin{itemize}
\item Step 0 — Verify signal quality
  \begin{itemize}
  \item Determine whether missing logs are present; missing logs are treated as \textbf{MISSING}/\textbf{UNKNOWN}, not inferred.
  \item Compliance thresholds govern any macro change; if thresholds unmet, no structural change.
  \end{itemize}
\item Step 1 — Lowest-cost correction
  \begin{itemize}
  \item Portable baseline concept: repeatable “default meal structure” posture with minimal decision points; maintain protein anchor and calorie posture.
  \end{itemize}
\item Step 2 — Next correction
  \begin{itemize}
  \item Hydration continuity per 10.C: baseline + session top-ups; electrolytes only if triggered; no novelty.
  \end{itemize}
\item Step 3 — Macro adjustment
  \begin{itemize}
  \item \textbf{HOLD} macro changes during travel unless safety dictates; if travel causes repeated under-eating and fatigue, use maintenance reset 7 days.
  \end{itemize}
\item Step 4 — Escalation, hold, refer
  \begin{itemize}
  \item Reslot gate tests if travel compromises validity or protected windows; do not compress make-ups per Stage 3.
  \item Hold Period Before Next Change: 7 days post-travel stabilization.
  \item What to Log: travel flag, compliance counts, hydration estimate, sleep markers, any reslots.
  \item Worked Example:
    \begin{itemize}
    \item Step 0: travel week missing logs 3 days; compliance cannot be verified.
    \item Step 1: portable baseline for remaining days; no macro changes.
    \item Step 4: reslot any planned test week impacted.
    \end{itemize}
  \end{itemize}
\end{itemize}

\newpage

\subsection*{LADDER-12 — Water Retention Perception Ladder — Education-Only Trend Discipline}
\phantomsection
\addcontentsline{toc}{subsection}{LADDER-12 — Water Retention Perception Ladder — Education-Only Trend Discipline}

Primary Signal: perceived water retention, scale volatility, frustration-driven urge to cut aggressively.\\
Confirm With Data: weekly average trend (\textbf{MET-2B-001}) across 7–14 days; sodium/carb swings; sleep and fatigue markers.\\
Stop Conditions Safety: restrictive behavior escalation; dehydration tactics attempted.\\
Review Cadence: Weekly.\\
Do Not Do: reactive cutting; dehydration practices; extreme sodium restriction; chugging/overhydration.\\
Steps in order:

\begin{itemize}
\item Step 0 — Verify signal quality
  \begin{itemize}
  \item Trend window: 7–14 days; decisions require weekly average trend, not single weigh-ins.
  \item Compliance threshold: structural changes only if $\geq$12/14 on-plan days.
  \end{itemize}
\item Step 1 — Lowest-cost correction
  \begin{itemize}
  \item Stabilize inputs: consistent sodium/carb patterns; stable meal timing; stable hydration routine per 10.C.
  \end{itemize}
\item Step 2 — Next correction
  \begin{itemize}
  \item Verify no novelty near gates; remove novelty foods that change GI and water retention perception.
  \end{itemize}
\item Step 3 — Macro adjustment
  \begin{itemize}
  \item Only if 14-day trend confirms true upward drift under high compliance: apply one calorie step (-150 to -250) and hold 7–14 days.
  \end{itemize}
\item Step 4 — Escalation, hold, refer
  \begin{itemize}
  \item If restrictive behavior distress escalates: maintenance reset 7 days; seek professional support posture; do not intensify restriction.
  \item Hold Period Before Next Change: 7–14 days after Step 3.
  \item What to Log: trend window, novelty events, compliance counts, hydration stability, action taken.
  \item Worked Example:
    \begin{itemize}
    \item Step 0: scale up 4 lb after salty meal; weekly average unchanged.
    \item Step 1: stabilize sodium/carb for 7 days; no cuts.
    \item Step 3: not applied because trend not confirmed.
    \end{itemize}
  \end{itemize}
\end{itemize}

\newpage

\section*{10.D.3 — Quick Reference Decision Tables}
\phantomsection
\addcontentsline{toc}{section}{10.D.3 — Quick Reference Decision Tables}

\subsection*{Symptom/Signal to Default Action}
\phantomsection
\addcontentsline{toc}{subsection}{Symptom/Signal to Default Action}

\begingroup
\scriptsize
\setlength{\tabcolsep}{2pt}
\renewcommand{\arraystretch}{1.12}

\begin{longtblr}{
  width=\linewidth,
  rowhead=1,
  colspec = {
    X[l,3.45]
    X[l,2.45]
    X[l,3.10]
  },
  hlines,
  vlines,
  cells = {hsep=6pt, vsep=6pt, halign=l, valign=t},
  row{odd} = {bg=white},
  row{even} = {bg=LightGray},
  row{1} = {bg=StageBlue!15, fg=black, font=\bfseries, halign=l, valign=m},
}
\Hdr{Signal} &
\Hdr{Default Action} &
\Hdr{Notes} \\
Chest pain/pressure, syncope/near-syncope, confusion, severe dizziness &
Stop and seek evaluation per governance &
Red-flag symptom stop posture overrides all ladders \\
Suspected heat illness signs &
Stop; cool down; seek evaluation; reslot tests &
Apply 10.C heat triggers; Stage 3 reslot discipline \\
Persistent vomiting/diarrhea &
Stop; seek evaluation; hold nutrition changes &
Hydration safety is primary; no forced drinking \\
Illness flags/systemic symptoms &
Hold deficit; maintenance reset if needed; reslot tests &
Do not test through illness; validity protection \\
\textgreater0.75\% \textbf{BW}/week loss $\times$ 2 weeks &
Increase calories +150–250; hold &
Protect recovery and validity; no novelty \\
Plateau definition met &
Simplify; then -150–250 if compliant; hold &
Requires $\geq$12/14 compliance and waist method consistency \\
Poor sleep trend &
Hold deficit; maintenance reset if collapsing &
No caffeine pattern change within 14 days of gates \\
Cramping/dehydration symptoms &
Apply 10.C sip model and electrolyte triggers &
Do not exceed 32 oz/hr; no chugging; no novelty products \\
Deload/Test week &
Calories to maintenance; no novelty &
Test-day carb redistribution allowed per 10.B \\
Gate proximity within 7–10 days &
Hold major changes &
Stability is the control; log any unavoidable change \\
\end{longtblr}

\endgroup

\subsection*{Problem Type to Allowed Change Type by Week Type}
\phantomsection
\addcontentsline{toc}{subsection}{Problem Type to Allowed Change Type by Week Type}

\begingroup
\scriptsize
\setlength{\tabcolsep}{2pt}
\renewcommand{\arraystretch}{1.12}

\begin{longtblr}{
  width=\linewidth,
  rowhead=1,
  colspec = {
    Q[l,wd=1.30in]
    X[l,3.05]
    X[l,2.65]
    X[l,2.80]
  },
  hlines,
  vlines,
  cells = {hsep=6pt, vsep=6pt, halign=l, valign=t},
  row{odd} = {bg=white},
  row{even} = {bg=LightGray},
  row{1} = {bg=StageBlue!15, fg=black, font=\bfseries, halign=l, valign=m},
}
\Hdr{Problem} &
\Hdr{Normal Training Week} &
\Hdr{Deload/Test Week} &
\Hdr{Notes} \\
Plateau &
Compliance simplification; then -150–250 if $\geq$12/14 &
No structural cuts; maintenance posture applies &
Reslot if validity would be contaminated \\
Too-fast loss &
+150–250 or maintenance reset if needed &
Maintenance posture applies &
Protect performance evidence \\
Hunger/cravings &
Simplify + distribution tweaks; +150 if persistent &
Maintenance posture; no novelty &
Avoid punitive restriction \\
Low energy/high RPE &
Maintenance reset 7 days &
Maintenance posture mandatory &
Do not cut during recovery collapse \\
GI distress &
Remove novelty; stabilize; hold &
No novelty; stabilize &
Reslot tests if symptoms invalidate validity \\
Heat/cold hazards &
Substitute/reslot; hydration per 10.C &
Reslot tests; indoor substitution &
Safety and validity override schedule \\
\end{longtblr}

\endgroup

\newpage

\section*{10.D.4 — Documentation Templates}
\phantomsection
\addcontentsline{toc}{section}{10.D.4 — Documentation Templates}

Troubleshooting Action Log template:

\begingroup
\scriptsize
\setlength{\tabcolsep}{2pt}
\renewcommand{\arraystretch}{1.12}

\begin{longtblr}{
  width=\linewidth,
  rowhead=1,
  colspec = {
    Q[l,wd=0.35in]
    Q[l,wd=0.55in]
    Q[l,wd=.50in]
    Q[l,wd=.50in]
    Q[l,wd=.65in]
    X[l,wd=.65in]
    Q[l,wd=0.5in]
    Q[l,wd=0.5in]
    Q[l,wd=0.5in]
    Q[l,wd=0.5in]
  },
  hlines,
  vlines,
  cells = {hsep=6pt, vsep=6pt, halign=l, valign=t},
  row{odd} = {bg=white},
  row{even} = {bg=LightGray},
  row{1} = {bg=StageBlue!15, fg=black, font=\bfseries, halign=l, valign=m},
}
\Hdr{Date} &
\Hdr{Phase/Week} &
\Hdr{Problem} &
\Hdr{Ladder Used} &
\Hdr{Step Executed} &
\Hdr{Change Made} &
\Hdr{Hold Period} &
\Hdr{Review Date} &
\Hdr{Outcome} &
\Hdr{Notes} \\
\end{longtblr}

\endgroup

Rules:

\begin{itemize}
\item Change Made must reflect 10.B levers:
  \begin{itemize}
  \item compliance fix,
  \item distribution tweak,
  \item $\pm$150–250 kcal/day,
  \item maintenance reset,
  \item hold and reslot.
  \end{itemize}
\item Notes must include Gate proximity field:
  \begin{itemize}
  \item within 14 days of gate windows Y/N,
  \item if Y, confirm no novelty posture applied.
  \end{itemize}
\end{itemize}

\newpage

\section*{10.D.5 — Integrity Checks}
\phantomsection
\addcontentsline{toc}{section}{10.D.5 — Integrity Checks}

\begin{itemize}
\item Trend-based verification precedes action; no single-day decisions.
\item Compliance/system fixes precede macro changes; thresholds ($\geq$6/7 and $\geq$12/14) gate changes.
\item Calorie changes use small-step control ($\pm$150–250) with 7–14 day hold periods.
\item Maintenance reset is exactly 7 days at maintenance and is used when recovery collapse or restrictive behavior risk appears.
\item No extreme restriction and no dehydration practices; water cuts are prohibited.
\item No novelty near gates/test weeks is enforced, including caffeine pattern change prohibition within 14 days of gate windows.
\item Deload/test weeks use maintenance-calorie posture and allow only test-day carb redistribution per 10.B.
\item Logging is required for every action using the Troubleshooting Action Log template; missing data is recorded as \textbf{MISSING}/\textbf{UNKNOWN}.
\item Escalation triggers and stop conditions are explicit and non-medical; red flags override ladders.
\item Ladders remain feasible under start-from-zero constraints and Stage 7.13 tooling posture.
\item Hydration references do not invent new targets; they use 10.C bands, triggers, ceilings, and reslot rules.
\end{itemize}

\end{document}